% Part: incompleteness
% Chapter: representability-in-q
% Section: representing-relations

\documentclass[../../../include/open-logic-section]{subfiles}

\begin{document}

\olfileid[id]{inc}{req}{rel}

\olsection{Merepresentasikan Relasi}

Mari kita nyatakan apa artinya suatu \emph{relasi} dapat direpresentasikan.

\begin{defn}
\ollabel{defn:representing-relations} Relasi $R(x_0,\dots,x_k)$ pada
bilangan asli {\em dapat direpresentasikan dalam $\Th{Q}$} jika terdapat
formula $!A_R(x_0,\dots,x_k)$ sedemikian sehingga, setiap kali
$R(n_0,\dots,n_k)$ benar, $\Th{Q}$ membuktikan
$!A_R(\num{n_0},\dots,\num{n_k})$, dan setiap kali
$R(n_0,\dots,n_k)$ salah, $\Th{Q}$ membuktikan
$\lnot !A_R(\num{n_0}, \dots, \num{n_k})$.
\end{defn}

\begin{thm}
\ollabel{thm:representing-rels} Suatu relasi dapat direpresentasikan dalam
$\Th{Q}$ jika dan hanya jika relasi itu dapat dihitung.
\end{thm}

\begin{proof}
Untuk arah maju, andaikan $R(x_0,\dots,x_k)$ direpresentasikan oleh formula
$!A_R(x_0,\dots,x_k)$. Berikut algoritme untuk menghitung $R$: dengan masukan
$n_0$, \dots,~$n_k$, carilah secara serentak pembuktian atas
$!A_R(\num{n_0}, \dots, \num{n_k})$ dan pembuktian atas
$\lnot !A_R(\num{n_0}, \dots, \num{n_k})$. Menurut hipotesis kita,
pencarian pasti menemukan salah satunya; jika yang ditemukan ialah yang
pertama, laporkan ``ya,'' dan jika tidak, laporkan ``tidak.''

Untuk arah lainnya, andaikan $R(x_0, \dots, x_k)$ dapat dihitung. Menurut
definisi, ini berarti bahwa fungsi $\Char{R}(x_0, \dots, x_k)$ dapat
dihitung. Menurut \olref[int]{thm:representable-iff-comp}, $\Char{R}$
direpresentasikan oleh suatu formula, misalnya
$!A_{\Char{R}}(x_0, \dots, x_k, y)$. Misalkan
$!A_R(x_0, \dots, x_k)$ merupakan formula
$!A_{\Char{R}}(x_0, \dots, x_k, \num{1})$. Maka, bagi sebarang
$n_0$, \dots,~$n_k$, jika $R(n_0, \dots, n_k)$ benar, maka
$\Char{R}(n_0, \dots, n_k) = 1$; dalam hal ini, $\Th{Q}$ membuktikan
$!A_{\Char{R}}(\num{n_0}, \dots, \num{n_k}, \num{1})$, sehingga
$\Th{Q}$ membuktikan $!A_R(\num{n_0}, \dots, \num{n_k})$. Di sisi lain,
jika $R(n_0, \dots, n_k)$ salah, maka
$\Char{R}(n_0, \dots, n_k) = 0$. Ini berarti bahwa $\Th{Q}$ membuktikan
\[
\lforall[y][(!A_{\Char{R}}(\num{n_0}, \dots, \num{n_k}, y) \lif y =
  \num{0})].
\]
Karena $\Th{Q}$ membuktikan $\eq/[\num{0}][\num{1}]$, $\Th{Q}$ membuktikan
$\lnot !A_{\Char{R}}(\num{n_0}, \dots, \num{n_k}, \num{1})$, dan dengan
demikian membuktikan $\lnot !A_R(\num{n_0}, \dots, \num{n_k})$.
\end{proof}

\begin{prob}
Tunjukkan bahwa jika $R$ dapat direpresentasikan dalam~$\Th{Q}$, demikian
pula~$\Char{R}$.
\end{prob}

\end{document}
